% content/c01-wahrscheinlichkeit/de.tex — Kapitel 1: Wahrscheinlichkeit für Physiker
% (deutsche Leitfassung)
%
% Migration des Unterabschnitts „Wahrscheinlichkeit“ der ersten Vorlesung
% (latex/lectures/de/vl01.tex:93-216).  Wohin jede Gleichung, jedes Beispiel und jede
% Übung gegangen ist, steht in evidence/migrations/maps/vl01-map.json.
%
% Definitionen: Registry D53 bis D56 (Quelle: das Skript, vl01.tex, mit Zeilenangabe).
% Herleitungen: jede als „hergeleitet“ markierte Aussage steht mit derselben ID in
%   evidence/derivations/c01/DERIVATION.md.
% Rahmung: docs/plan/BRIEF.md:45-46 (author-stipulated); Gefahr H1.1 des Curriculum-Berichts
%   (Determinismus ist der Sonderfall des Punktgewichts) in Proposition c01.prop.punktgewicht.
%
% Stimme (Überarbeitung „Lecture 1“, 2026-10-20): Die Prosa folgt TJOs Sätzen aus vl01.tex
%   und seinen Vorlesungsvideos v01 (ab [60:37]) und v02 (bis [75:00]); Zeitmarken stehen als
%   Kommentar an der Stelle.  Die Transkripte sind keine registrierte Quelle; aus ihnen
%   kommen Wortlaut, Beispiele (zwei Münzen, das Rubbellos, das durchgerechnete Beispiel
%   mit drei Ergebnissen) und didaktische Bemerkungen.  Die Rechnungen der Beispiele sind
%   hergeleitet (DERIVATION.md, Teil B).
%
% Schreibweise (Registry): Ergebnismenge X = \Outcomes, Gewicht p = \pw, \Prob, \Ex, \Var,
% Indikatorfunktion \ind{S}.  Das Punktgewicht auf a heißt hier p_a, das Produktgewicht
% für N Wiederholungen p^{(N)}; beide sind Schreibweisen dieses Kapitels.

\chapter{Wahrscheinlichkeit für Physiker}\label{c01}

% Stimme: vl01.tex:94; v01 [60:37]-[64:24] („ein intuitives Verständnis von
% Wahrscheinlichkeiten setze ich voraus“).
Die Zahlen, welche die Quantenmechanik als Ergebnis liefert, sind \emph{Wahrscheinlichkeiten}.
Wie sind solche Größen zu verstehen, und wie arbeiten wir damit? Der richtige Rahmen dafür ist
die Wahrscheinlichkeitstheorie, und die bauen wir in diesem Kapitel auf. Erwarten Sie nichts
konzeptionell Neues: Ein intuitives Verständnis von Wahrscheinlichkeiten, Verteilungen und
Erwartungswerten setze ich voraus. Es geht um die \emph{Notation}, in der wir später die
Quantenmechanik formulieren, und um eine Sprache. Dazu gehört eine Festlegung, die ich für
den ganzen Kurs treffe:
% GAP-TJO: Deutsche Fassung von docs/plan/BRIEF.md:46 (erster Teil des Satzes): Übersetzung
% des Kapitelautors (G-0022).  Der zweite Teil des Satzes (Messungen als Elemente des
% Dualraums) gehört zu Kapitel 2 und ist hier nicht zitiert.
\begin{equation}\label{c01.eq.sprache}\claim{c01.cl.sprache}{author-stipulated}
  \begin{aligned}
    &\text{\tpcifenglish{Thus our presentation starts with probability theory,}{Unsere Darstellung beginnt deshalb mit der}} \\
    &\text{\tpcifenglish{with notation and framing}{Wahrscheinlichkeitstheorie, mit Notation und Rahmung,}} \\
    &\text{\tpcifenglish{adapted to physics: we use the language}{die der Physik angepasst sind: Wir benutzen die Sprache}} \\
    &\text{\tpcifenglish{of preparations (= states)}{der Präparationen (= Zustände)}} \\
    &\text{\tpcifenglish{instead of “probability measures”.}{statt der „Wahrscheinlichkeitsmaße“.}}
  \end{aligned}
\end{equation}
Was die Mathematik eine Wahrscheinlichkeitsverteilung nennt, heißt hier also Präparation,
und was sie eine Zufallsvariable nennt, heißt Messung. Tabelle~\ref{c01.tab.sprache} am Ende
von Abschnitt~\ref{c01.s2} stellt die Wörter nebeneinander.

% Stimme: v02 [00:00] („alle drei müssen spezifiziert sein, bevor man überhaupt anfangen
% kann zu rechnen“).
Um ein Experiment wahrscheinlichkeitstheoretisch zu modellieren, brauchen wir genau drei
Dinge: die Menge der möglichen Ergebnisse, die Präparation und die Messung. Alle drei müssen
festgelegt sein, bevor wir überhaupt anfangen können zu rechnen. Wir nehmen sie in dieser
Reihenfolge.

\section{Ergebnisse und Präparationen}\label{c01.s1}

% Stimme: vl01.tex:98-101; v01 [64:24] („man erhält ihn, indem man alle möglichen
% Ergebnisse auflistet“; zwei Münzen an der Tafel).
Der erste Schritt bei der Modellierung jedes probabilistischen Experiments ist, ihm eine
Menge zuzuordnen. Man erhält sie, indem man alle möglichen Ergebnisse auflistet.

\tpcdef{D53}

In diesem und den nächsten beiden Abschnitten ist $\Outcomes$ eine endliche Menge.

\begin{example}[Würfel und Münzen]\label{c01.ex.wuerfel}
  Bei einem Würfel besteht die Ergebnismenge aus den möglichen Augenzahlen:
  $\Outcomes = \{1, 2, 3, 4, 5, 6\}$. Bei einer Münze ist
  $\Outcomes = \{\text{\tpcifenglish{heads}{Kopf}}, \text{\tpcifenglish{tails}{Zahl}}\}$,
  bei zwei Münzen mit K für Kopf und Z für Zahl
  $\Outcomes = \{\text{KK}, \text{KZ}, \text{ZK}, \text{ZZ}\}$. „Die Augenzahl ist gerade“
  ist das Ereignis $S = \{2, 4, 6\}$.
\end{example}

\begin{exercise}[Zwei Würfel]\label{c01.exc.zwei-wuerfel}\claim{c01.cl.zwei-wuerfel}{derived}
  Beschreiben Sie die Ergebnismenge für \emph{zwei} Würfel. Wie viele Elemente hat sie?
\end{exercise}
\begin{solution}
  Ein Ergebnis ist ein Paar $(x_1, x_2)$ aus der Augenzahl $x_1$ des ersten und der
  Augenzahl $x_2$ des zweiten Würfels. Die Ergebnismenge ist die Menge aller solchen Paare
  mit $x_1, x_2 \in \{1, \dots, 6\}$. Zu jedem der $6$ Werte von $x_1$ gibt es $6$ Werte
  von $x_2$; sie hat also $6 \cdot 6 = 36$ Elemente
  (Abbildung~\ref{c01.fig.ereignis}).
\end{solution}

\begin{tpcfigure}{c01.fig.ereignis}
  \caption{Ergebnismengen und Ereignisse. (a) Ein Würfel: sechs Ergebnisse; das Ereignis
    „gerade Augenzahl“ ist die Menge der schraffierten Felder. (b) Zwei Würfel: Jeder
    Punkt ist ein Paar $(x_1, x_2)$, zusammen $36$ Ergebnisse. Das Ereignis „die
    Augensumme ist $7$“ besteht aus den sechs mit Quadraten markierten Paaren.}
\end{tpcfigure}

% Stimme: vl01.tex:108-116; v01 [70:32] („wie modelliert man einen fairen Würfel, eine faire
% Münze oder einen gezinkten Würfel?“).
Schritt zwei ist die Präparation. Wie modelliert man einen fairen Würfel, eine faire Münze
oder einen gezinkten Würfel? In einer probabilistischen Theorie durch die Angabe einer
Wahrscheinlichkeit für jedes Ergebnis.

\tpcdef{D54}

Die beiden Bedingungen an $\pw$ sind die einzigen. Die erste gilt, weil eine negative
Wahrscheinlichkeit nicht interpretierbar ist; die zweite, weil die Wahrscheinlichkeit, dass
irgendein Ergebnis eintritt, immer eins ist. Für die faire Münze ist
$\pw(\text{\tpcifenglish{heads}{Kopf}}) = \pw(\text{\tpcifenglish{tails}{Zahl}}) = \tfrac12$.
Wir nennen $\pw$ im Folgenden kurz ein \emph{Gewicht}.

\begin{example}[Fairer und gezinkter Würfel]\label{c01.ex.gezinkt}\claim{c01.cl.gezinkt}{derived}
  Für einen fairen Würfel ist $\pw(x) = \tfrac16$ für alle $x$. Andere Gewichte beschreiben
  \emph{gezinkte} Würfel. In diesem Kapitel dient als gezinkter Würfel das Gewicht
  \begin{equation}\label{c01.eq.gezinkt}
    \pw(1) = \tfrac{1}{12}, \qquad
    \pw(2) = \pw(3) = \pw(4) = \pw(5) = \tfrac{2}{12}, \qquad
    \pw(6) = \tfrac{3}{12} .
  \end{equation}
  Die Summe ist $\tfrac{1 + 4\cdot 2 + 3}{12} = 1$. Das Ereignis „gerade Augenzahl“,
  $S = \{2,4,6\}$, hat beim fairen Würfel die Wahrscheinlichkeit
  $\Prob[S] = 3\cdot\tfrac16 = \tfrac12$ und beim gezinkten
  $\Prob[S] = \tfrac{2 + 2 + 3}{12} = \tfrac{7}{12}$.
\end{example}

\begin{tpcfigure}{c01.fig.gewichte}
  \caption{Drei Gewichte auf der Ergebnismenge $\{1, \dots, 6\}$ eines Würfels. Die Höhe
    eines Balkens ist $\pw(x)$; in jedem Bild ist die Summe aller Höhen $1$. (a) Fairer
    Würfel. (b) Der gezinkte Würfel aus Gleichung~\eqref{c01.eq.gezinkt}. (c) Ein
    Punktgewicht: Das ganze Gewicht liegt auf dem Ergebnis $3$. Solche Gewichte beschreiben
    deterministische Präparationen (Proposition~\ref{c01.prop.punktgewicht}).}
\end{tpcfigure}

\begin{proposition}[Eigenschaften der Wahrscheinlichkeit]\label{c01.prop.wahrscheinlichkeit}\claim[proof-included]{c01.cl.wahrscheinlichkeit}{derived}
  Sei $\pw$ ein Gewicht auf einer endlichen Ergebnismenge $\Outcomes$, und seien
  $S, T \subset \Outcomes$ Ereignisse. Dann gilt:
  \begin{enumerate}
    \item $0 \leq \Prob[S] \leq 1$, $\Prob[\emptyset] = 0$ und $\Prob[\Outcomes] = 1$.
    \item Sind $S_1, \dots, S_k$ paarweise disjunkt, das heißt, haben je zwei von ihnen
      kein Ergebnis gemeinsam, so ist
      \begin{equation}\label{c01.eq.additiv}
        \Prob[S_1 \cup \dots \cup S_k] = \Prob[S_1] + \dots + \Prob[S_k] .
      \end{equation}
    \item $\Prob[\Outcomes \setminus S] = 1 - \Prob[S]$.
    \item $\Prob[S \cup T] = \Prob[S] + \Prob[T] - \Prob[S \cap T]$.
  \end{enumerate}
\end{proposition}
\begin{proof}
  Zu~2: Jedes Ergebnis $x$ der Vereinigung liegt in genau einem $S_j$. Die Summe der
  $\pw(x)$ über die Vereinigung lässt sich deshalb nach den $S_j$ ordnen, und das ist die
  rechte Seite. Zu~1: $\Prob[S]$ ist eine Summe von Zahlen $\pw(x) \geq 0$, also
  $\Prob[S] \geq 0$; die leere Summe ist $0$, und $\Prob[\Outcomes] = 1$ ist die zweite
  Bedingung an ein Gewicht. Zu~3: $S$ und $\Outcomes \setminus S$ sind disjunkt mit
  Vereinigung $\Outcomes$, also ist nach~2 $\Prob[S] + \Prob[\Outcomes \setminus S] = 1$.
  Daraus folgt auch $\Prob[S] \leq 1$. Zu~4: $S \cup T$ ist die disjunkte Vereinigung von
  $S$ und $T \setminus S$, und $T$ ist die disjunkte Vereinigung von $T \setminus S$ und
  $S \cap T$. Nach~2 ist $\Prob[S \cup T] = \Prob[S] + \Prob[T \setminus S]$ und
  $\Prob[T] = \Prob[T \setminus S] + \Prob[S \cap T]$. Einsetzen ergibt die Behauptung.
\end{proof}

Die Disjunktheit in~\eqref{c01.eq.additiv} ist wesentlich. Überlappen sich die Ereignisse,
so zählt die rechte Seite Ergebnisse doppelt; Teil~4 zeigt für zwei Ereignisse, um wie viel.
Eine Funktion, die Mengen Zahlen $\geq 0$ zuordnet und \eqref{c01.eq.additiv} erfüllt, heißt in der
Mathematik ein \emph{Maß}, und ein \emph{Wahrscheinlichkeitsmaß}, wenn
$\Prob[\Outcomes] = 1$ ist. Mehr Maßtheorie braucht dieser Kurs im Kern nicht; die Wörter
sollen nur die Literatur lesbar machen.

Präparationen lassen sich mischen. Sind $\pw$ und $q$ Gewichte und ist $\lambda$ eine Zahl
zwischen $0$ und $1$, so kann man mit Wahrscheinlichkeit $\lambda$ nach $\pw$ und sonst
nach $q$ präparieren. Auch das ist eine Präparation, mit dem Gewicht
\begin{equation}\label{c01.eq.mischung}\claim{c01.cl.mischung}{derived}
  \lambda\, \pw + (1-\lambda)\, q,
  \qquad
  \bigl(\lambda\, \pw + (1-\lambda)\, q\bigr)(x) = \lambda\, \pw(x) + (1-\lambda)\, q(x) .
\end{equation}
Das ist wieder ein Gewicht: Alle Werte sind $\geq 0$, weil $\lambda$, $1-\lambda$,
$\pw(x)$ und $q(x)$ es sind, und die Summe über $x$ ist
$\lambda \cdot 1 + (1-\lambda)\cdot 1 = 1$.

\section{Messungen und Erwartungswert}\label{c01.s2}

% Stimme: vl01.tex:118-126; v01 [77:53]; v02 [04:15] („das Experiment selbst bildet
% Wahrscheinlichkeitsfunktionen auf Wahrscheinlichkeitsfunktionen ab; genau das wird in der
% Quantenmechanik immer wieder vorkommen“).
Kapitel~0 hat jedes Experiment in drei Stadien zerlegt. In der klassischen
Wahrscheinlichkeitstheorie hat jedes Stadium eine einfache mathematische Gestalt
\srcref{tpc-script-v0.2}{vl01, Kasten „Jedes Experiment in drei Schritten“}.

\begin{greybox}[Die drei Stadien, klassisch]
  \begin{equation}\label{c01.eq.stadien}\claim{c01.cl.messung}{source-backed}
    \begin{aligned}
      &\text{\tpcifenglish{1. Preparation:}{1. Präparation:}} && \text{\tpcifenglish{a weight}{ein Gewicht}}\ \pw \\
      &\text{\tpcifenglish{2. Run:}{2. Durchführung:}} && \text{\tpcifenglish{a map of weights to weights}{eine Abbildung von Gewichten auf Gewichte}} \\
      &\text{\tpcifenglish{3. Measurement:}{3. Messung:}} && \text{\tpcifenglish{a function}{eine Funktion}}\ f\colon \Outcomes \to \mathbb{R}
    \end{aligned}
  \end{equation}
  Am Ende der Messung steht eines der möglichen Ergebnisse. Was für Abbildungen im zweiten
  Stadium vorkommen, braucht dieses Kapitel nicht; in der Quantenmechanik wird uns dieses
  Bild immer wieder begegnen.
\end{greybox}

\tpcdef{D55}

% Stimme: vl01.tex:133; v01 [81:48]-[83:27] („quasi alle Funktionen dürfen als Messung
% interpretiert werden“), [86:09].
Beachten Sie, wie großzügig das ist: An eine Messung stellen wir keine Bedingung, quasi jede
Funktion auf $\Outcomes$ darf als Messung gelten. An die Präparation stellen wir dagegen
zwei. Der Wert $f(x)$ ist der Messwert, der sich ergibt, wenn man sicher weiß, dass das
Ergebnis $x$ vorliegt. Der Erwartungswert verbindet Präparation und Messung: Er ist die mit
den Wahrscheinlichkeiten gewichtete Summe der Messwerte.

\begin{example}[Augenzahl]\label{c01.ex.augenzahl}\claim{c01.cl.augenzahl}{derived}
  Für einen Würfel ist die Augenzahl $f(x) = x$ eine Messung. Für den fairen Würfel ist
  \begin{equation}\label{c01.eq.augenzahl-fair}
    \Ex[f] = \tfrac16\,(1 + 2 + 3 + 4 + 5 + 6) = \tfrac{21}{6} = \tfrac72 ,
  \end{equation}
  für den gezinkten Würfel aus Gleichung~\eqref{c01.eq.gezinkt}
  \begin{equation}\label{c01.eq.augenzahl-gezinkt}
    \Ex[f] = \tfrac{1}{12}\,(1\cdot 1 + 2\cdot 2 + 2\cdot 3 + 2\cdot 4 + 2\cdot 5 + 3\cdot 6)
           = \tfrac{47}{12} .
  \end{equation}
  Keiner der beiden Werte ist eine Augenzahl. Der Erwartungswert muss kein möglicher
  Messwert sein. Die Augenzahl ist auch nicht die einzige Messung am Würfel: „Ist die
  Augenzahl gerade?“ und „Ist die Augenzahl gleich sechs?“ sind ebenfalls Messungen, mit
  Ja oder Nein als Antwort.
\end{example}

\begin{tpcfigure}{c01.fig.erwartung}
  \caption{Der Erwartungswert der Augenzahl als Schwerpunkt. Denkt man sich die Gewichte
    $\pw(x)$ als Massen an den Stellen $x$ eines Balkens, so ist $\Ex[f]$ die Stelle, an
    der der Balken im Gleichgewicht ist: (a) beim fairen Würfel in der Mitte, bei
    $\tfrac72$; (b) beim gezinkten Würfel aus Gleichung~\eqref{c01.eq.gezinkt} nach rechts
    verschoben, bei $\tfrac{47}{12}$.}
\end{tpcfigure}

\subsection{Relative Häufigkeiten}\label{c01.s2.1}

% Stimme: vl01.tex:137; v02 [06:44] (der Erwartungswert wird nicht vom Himmel geworfen,
% sondern motiviert).
Woher kommt die Formel für den Erwartungswert? Sie fällt nicht vom Himmel. Man wiederholt
das Experiment $N$-mal mit
derselben Präparation, wobei $N$ eine positive ganze Zahl ist, und notiert die Ergebnisse
$x_1, \dots, x_N$. Sei $N(x)$ die Anzahl
der Durchläufe mit dem Ergebnis $x$. Die \emph{relative Häufigkeit} von $x$ ist
\begin{equation}\label{c01.eq.haeufigkeit}\claim{c01.cl.haeufigkeit}{derived}
  h_N(x) = \frac{N(x)}{N},
  \qquad
  h_N(x) \geq 0,
  \qquad
  \sum_{x \in \Outcomes} h_N(x) = 1 .
\end{equation}
Die Summe ist $1$, weil jeder der $N$ Durchläufe genau ein Ergebnis hat, also
$\sum_x N(x) = N$ ist. Die relativen Häufigkeiten einer Messreihe sind selbst ein Gewicht.
Für den Mittelwert der Messwerte einer Messung $f$ über die Messreihe gilt
\begin{equation}\label{c01.eq.stichprobenmittel}\claim{c01.cl.stichprobenmittel}{derived}
  \frac{1}{N} \sum_{i=1}^{N} f(x_i)
  = \frac{1}{N} \sum_{x \in \Outcomes} N(x)\, f(x)
  = \sum_{x \in \Outcomes} h_N(x)\, f(x) ,
\end{equation}
denn in der linken Summe kommt der Summand $f(x)$ genau $N(x)$-mal vor. Der Mittelwert
einer Messreihe ist also der Erwartungswert von $f$ im Gewicht $h_N$. Diese Gleichung ist
exakt und gilt für jede Messreihe.

Und wie hängt $h_N$ mit $\pw$ zusammen? Führen wir das Experiment $N$-mal unabhängig mit
derselben Präparation durch, so liegt bei vielen Wiederholungen die relative Häufigkeit
$h_N(x)$ mit hoher Wahrscheinlichkeit nahe bei $\pw(x)$,
\begin{equation}\label{c01.eq.lesart}\claim{c01.cl.lesart}{source-backed}
  h_N(x) \approx \pw(x)
  \qquad \text{\tpcifenglish{for large}{für große}}\ N .
\end{equation}
In diesem Modell kommt das Ergebnis $x$ bei großem $N$ mit hoher Wahrscheinlichkeit
etwa $N \pw(x)$-mal vor, die Summe der Messwerte ist dann etwa
$\sum_x N \pw(x) f(x)$, und nach Division durch $N$ steht der Erwartungswert da
\srcref{tpc-script-v0.2}{vl01, vor der Formel für den Erwartungswert}. Diese Lesart ist
eine Deutung im angegebenen Modell: In einer einzelnen Messreihe kann $h_N$ von $\pw$
abweichen. Für unabhängige Wiederholungen derselben Präparation lässt sich eine obere
Schranke für die Wahrscheinlichkeit großer Abweichungen beweisen\zoomref{c01.v2}.

\begin{tpcfigure}{c01.fig.haeufigkeit}
  \demo{c01.demo.haeufigkeit}
  \caption{Eine Messreihe mit dem gezinkten Würfel aus Gleichung~\eqref{c01.eq.gezinkt},
    erzeugt von einem Rechenprogramm mit festem Startwert. Schraffierte Balken: relative
    Häufigkeiten $h_N(x)$ nach den ersten (a) $12$, (b) $120$ und (c) $1200$ Würfen
    derselben Reihe. Dicke Querstriche: die Gewichte $\pw(x)$. In (a) liegen die Balken
    weit von den Strichen, in (c) nah an ihnen. Die Balken eines Bildes haben zusammen
    stets die Höhe $1$, wie in Gleichung~\eqref{c01.eq.haeufigkeit}.}
\end{tpcfigure}

\subsection{Rechnen mit Erwartungswerten}\label{c01.s2.2}

Messungen sind Funktionen und lassen sich deshalb addieren, mit Zahlen multiplizieren und
miteinander multiplizieren: $(\alpha f + \gamma g)(x) = \alpha f(x) + \gamma g(x)$ und
$(fg)(x) = f(x)\,g(x)$. Eine Zahl $c$ wird dabei als die konstante Messung mit dem Wert $c$
gelesen. Wir schreiben $f \leq g$, wenn $f(x) \leq g(x)$ für alle $x \in \Outcomes$ gilt.

\begin{proposition}[Rechenregeln für den Erwartungswert]\label{c01.prop.erwartung}\claim[proof-included]{c01.cl.erwartung}{derived}
  Sei $\pw$ ein Gewicht auf einer endlichen Ergebnismenge $\Outcomes$, seien $f$ und $g$
  Messungen und $\alpha$, $\gamma$, $c$ reelle Zahlen. Dann gilt:
  \begin{enumerate}
    \item Linearität: $\Ex[\alpha f + \gamma g] = \alpha\, \Ex[f] + \gamma\, \Ex[g]$.
    \item Normierung: $\Ex[c] = c$.
    \item Monotonie: Aus $f \leq g$ folgt $\Ex[f] \leq \Ex[g]$. Insbesondere ist
      $\Ex[f] \geq 0$, wenn $f \geq 0$ ist.
    \item Der Erwartungswert liegt zwischen dem kleinsten und dem größten Messwert:
      $\min_x f(x) \leq \Ex[f] \leq \max_x f(x)$.
  \end{enumerate}
\end{proposition}
\begin{proof}
  Zu~1:
  \begin{equation}\label{c01.eq.linear}
    \sum_{x \in \Outcomes} \pw(x)\,\bigl(\alpha f(x) + \gamma g(x)\bigr)
    = \alpha \sum_{x \in \Outcomes} \pw(x) f(x) + \gamma \sum_{x \in \Outcomes} \pw(x) g(x) .
  \end{equation}
  Zu~2: $\sum_x \pw(x)\, c = c \sum_x \pw(x) = c$. Zu~3: Nach~1 ist
  $\Ex[g] - \Ex[f] = \Ex[g - f] = \sum_x \pw(x)\,(g(x) - f(x))$, und jeder Summand ist ein
  Produkt zweier Zahlen $\geq 0$. Zu~4: Mit $m = \min_x f(x)$ und $M = \max_x f(x)$ ist
  $m \leq f \leq M$; nach~3 und~2 folgt $m \leq \Ex[f] \leq M$.
\end{proof}

\subsection{Ja-Nein-Messungen}\label{c01.s2.3}

% Stimme: vl01.tex:142; v01 [86:09] („die meisten Messungen in der Quantenmechanik sind von
% dieser Art: Teilchen ist da, Teilchen ist nicht da“); v02 [12:00]-[25:40].
Von den unendlich vielen möglichen Messungen sind die wichtigsten die, die nur die Werte $0$
und $1$ annehmen. Sie sind das Modell für Detektoren, die das Vorhandensein eines Teilchens
anzeigen, etwa einen Photonendetektor: Der Detektor sagt Klick oder kein Klick. An diesen
zweiwertigen Messungen hängt die Quantenmechanik.
Eine solche Messung $f$ ist durch die Menge $S = \{x \in \Outcomes \colon f(x) = 1\}$
bestimmt. Man schreibt
\begin{equation}\label{c01.eq.indikator-def}\claim{c01.cl.indikator-def}{source-backed}
  \ind{S}(x) =
  \begin{cases}
    1, & x \in S, \\
    0, & \text{\tpcifenglish{otherwise}{sonst}}
  \end{cases}
\end{equation}
und nennt $\ind{S}$ die \emph{Indikatorfunktion} des Ereignisses $S$. Funktion und Menge
enthalten genau dieselbe Information. Halten Sie fest: Indikatorfunktionen \emph{sind}
Messungen, kein eigenes Objekt daneben. Die Menge $S$ darf beliebig sein. Drei Grenzfälle:
$S = \Outcomes$ gibt die konstante Messung $1$, $S = \emptyset$ die konstante Messung $0$,
und für ein einzelnes Ergebnis $a$ ist $\ind{\{a\}}(x)$ gleich $1$ für $x = a$ und sonst
$0$; das ist das Kronecker-Delta.

% Stimme: v02 [17:24] (das Rubbellos).
\begin{remark}
  Indikatorfunktionen begegnen einem überall, wo Wahrscheinlichkeiten vorkommen. Beim
  Glücksspiel ist die wichtigste Messung überhaupt: Habe ich gewonnen oder nicht? Bei einem
  Rubbellos ist die Ergebnismenge die Menge aller Zahlenkombinationen, die auftauchen
  können. Die Messung ist $1$ bei Gewinn und $0$ sonst. Ausgeschrieben ist diese Funktion
  enorm kompliziert, denn man muss jede Kombination daraufhin prüfen, ob sie gewinnt; und
  doch ist sie nur eine Indikatorfunktion. Die ganze Komplexität steckt in der Menge $S$.
\end{remark}

Der Erwartungswert einer Indikatorfunktion ist die Wahrscheinlichkeit ihres Ereignisses:
\begin{equation}\label{c01.eq.indikator}\claim{c01.cl.indikator}{derived}
  \Ex[\ind{S}] = \sum_{x \in \Outcomes} \pw(x)\, \ind{S}(x) = \sum_{x \in S} \pw(x) = \Prob[S] .
\end{equation}
Abbildung~\ref{c01.fig.indikator} zeigt die Rechnung für den gezinkten Würfel. Der
Erwartungswert einer Indikatorfunktion hat also eine Sonderform: Er ist eine reine Summe von
Wahrscheinlichkeiten. Wahrscheinlichkeiten sind besondere Erwartungswerte. Umgekehrt legt $\Prob$ alle
Erwartungswerte fest, denn mit $S = \{x\}$ ist $\pw(x) = \Prob[\{x\}]$ und deshalb
$\Ex[f] = \sum_{x} \Prob[\{x\}]\, f(x)$.

\begin{tpcfigure}{c01.fig.indikator}
  \caption{Eine Ja-Nein-Messung am gezinkten Würfel aus Gleichung~\eqref{c01.eq.gezinkt}.
    (a) Die Indikatorfunktion des Ereignisses $S = \{2,4,6\}$, gestrichelt wie alle
    Messungen. (b) Die Gewichte $\pw(x)$ (Umrisse) und die Produkte $\pw(x)\,\ind{S}(x)$
    (schraffiert): Die Indikatorfunktion lässt nur die Balken über $S$ übrig. Ihre Höhen
    ergeben zusammen $\Ex[\ind{S}] = \Prob[S] = \tfrac{7}{12}$.}
\end{tpcfigure}

Mit Indikatorfunktionen rechnet man wie mit den Mengen:
\begin{equation}\label{c01.eq.indikator-regeln}\claim{c01.cl.indikator-regeln}{derived}
  \ind{S}\,\ind{T} = \ind{S \cap T},
  \qquad
  \ind{S}^2 = \ind{S},
  \qquad
  1 - \ind{S} = \ind{\Outcomes \setminus S} .
\end{equation}
Das prüft man Ergebnis für Ergebnis: Das Produkt $\ind{S}(x)\,\ind{T}(x)$ ist genau dann
$1$, wenn beide Faktoren $1$ sind, also wenn $x$ in $S$ und in $T$ liegt. Die zweite
Gleichung ist der Fall $T = S$. Die dritte: $1 - \ind{S}(x)$ ist $0$ für $x \in S$ und $1$
sonst.

Jede Messung lässt sich aus Ja-Nein-Messungen zusammensetzen. Für einen Messwert $v$ sei
$\{f = v\}$ das Ereignis $\{x \in \Outcomes \colon f(x) = v\}$, „die Messung ergibt $v$“.
Dann ist
\begin{equation}\label{c01.eq.werteverteilung}\claim{c01.cl.werteverteilung}{derived}
  f = \sum_{v} v\, \ind{\{f = v\}},
  \qquad
  \Ex[f] = \sum_{v} v\, \Prob[\{f = v\}] ,
\end{equation}
wobei $v$ die verschiedenen Werte durchläuft, die $f$ annimmt. Die erste Gleichung gilt,
weil an jeder Stelle $x$ genau eine Indikatorfunktion gleich eins ist, die mit
$v = f(x)$. Ist $f(x) = 0$, verschwinden alle gewichteten Summanden. Die zweite folgt daraus mit der Linearität und Gleichung~\eqref{c01.eq.indikator}. Der
Erwartungswert ist das mit den Wahrscheinlichkeiten gewichtete Mittel der Messwerte.

% Stimme: v02 [38:07]-[48:44] (das durchgerechnete Beispiel der Vorlesung: Photonenzahl,
% alle 2^3 Teilmengen, „wie groß ist die Wahrscheinlichkeit, dass kein Ereignis vorliegt?
% Es liegt immer eines vor“, „wir haben längst Erwartungswerte berechnet“, negative
% Messwerte, Drehimpuls).
Ein vollständig durchgerechnetes Beispiel zeigt den ganzen Apparat im Kleinen. In der
Quantenmechanik werden wir sehr viele Erwartungswerte berechnen; das soll Ihnen am Ende sehr
leicht fallen. Es lohnt sich, viele solche Rechnungen selbst zu machen.

\begin{example}[Ein Experiment mit drei Ergebnissen]\label{c01.ex.drei-ergebnisse}\claim{c01.cl.drei-ergebnisse}{derived}
  Die drei Ergebnisse von $\Outcomes = \{1, 2, 3\}$ mögen etwa die Anzahl der Photonen sein,
  die ein Detektor zählt. Als Präparation wählen wir, beliebig, aber zulässig,
  $\pw(1) = \tfrac12$ und $\pw(2) = \pw(3) = \tfrac14$: Alle Werte sind $\geq 0$, und
  $\tfrac12 + \tfrac14 + \tfrac14 = 1$. Eine Menge mit drei Elementen hat $2^3 = 8$
  Teilmengen; jede ist ein Ereignis und hat ihre Indikatorfunktion.
  Tabelle~\ref{c01.tab.ereignisse} listet alle acht mit ihren Wahrscheinlichkeiten. Wie groß
  ist die Wahrscheinlichkeit, dass kein Ergebnis vorliegt? Es liegt immer eines vor, also
  ist $\Prob[\emptyset] = 0$. Die Einermengen haben die Werte von $\pw$, die Zweiermengen
  folgen aus der Additivität~\eqref{c01.eq.additiv}, und $\Prob[\Outcomes] = 1$. Damit ist
  die gesamte Information über das Experiment aufgelistet, und nebenbei haben wir schon acht
  Erwartungswerte berechnet: $\Prob[\{2,3\}] = \tfrac12$ \emph{ist} der Erwartungswert der
  Indikatorfunktion $\ind{\{2,3\}}$.

  Nun eine Messung, die keine Indikatorfunktion ist: $f(1) = -1$, $f(2) = 0$, $f(3) = +1$.
  Ihr Erwartungswert ist
  \begin{equation}\label{c01.eq.drei-ergebnisse}
    \Ex[f] = \tfrac12\cdot(-1) + \tfrac14\cdot 0 + \tfrac14\cdot(+1) = -\tfrac14 .
  \end{equation}
  Er ist negativ, und das ist kein Problem. Nur Wahrscheinlichkeiten müssen $\geq 0$ sein;
  Messwerte und Erwartungswerte dürfen negativ sein.
\end{example}

\begin{tpctable}{c01.tab.ereignisse}
  \caption{Alle acht Ereignisse des Experiments aus Beispiel~\ref{c01.ex.drei-ergebnisse},
    mit den Werten ihrer Indikatorfunktion an den Stellen $1$, $2$, $3$ und ihrer
    Wahrscheinlichkeit für $\pw = (\tfrac12, \tfrac14, \tfrac14)$. Jede Zeile ist eine
    Ja-Nein-Messung und ihr Erwartungswert.}
  \begin{tabular}{lcccc}
    \toprule
    $S$ & $\ind{S}(1)$ & $\ind{S}(2)$ & $\ind{S}(3)$ & $\Prob[S] = \Ex[\ind{S}]$ \\
    \midrule
    $\emptyset$ & $0$ & $0$ & $0$ & $0$ \\
    $\{1\}$ & $1$ & $0$ & $0$ & $1/2$ \\
    $\{2\}$ & $0$ & $1$ & $0$ & $1/4$ \\
    $\{3\}$ & $0$ & $0$ & $1$ & $1/4$ \\
    $\{1, 2\}$ & $1$ & $1$ & $0$ & $3/4$ \\
    $\{1, 3\}$ & $1$ & $0$ & $1$ & $3/4$ \\
    $\{2, 3\}$ & $0$ & $1$ & $1$ & $1/2$ \\
    $\{1, 2, 3\}$ & $1$ & $1$ & $1$ & $1$ \\
    \bottomrule
  \end{tabular}
\end{tpctable}

Messwerte $-1$, $0$, $+1$ sind keine Erfindung für die Übung. Beim Drehimpuls werden wir
in Kapitel~10 Zustände $\ket{\lqn, \mqn}$ kennenlernen, auf denen die
Drehimpulskomponente $\Lj{3}$ den Wert $\hbar \mqn$ hat, mit
\begin{equation}\label{c01.eq.drehimpuls}\claim{c01.cl.drehimpuls}{source-backed}
  \mqn \in \{-\lqn, -\lqn + 1, \dots, \lqn\},
  \qquad
  \lqn = 1\colon \quad \mqn \in \{-1, 0, 1\} .
\end{equation}
In Einheiten von $\hbar$ sind das für $\lqn = 1$ genau die Messwerte von $f$
\srcref{tpc-script-v0.2}{vl14, Eigenwerte des Drehimpulses}.

\begin{exercise}[Drei Ergebnisse]\label{c01.exc.drei}\claim{c01.cl.drei}{derived}
  Sei $\Outcomes = \{1, 2, 3\}$ und $\pw(x) = \tfrac13$ für alle $x \in \Outcomes$.
  Berechnen Sie $\Prob[S]$ für $S = \{1, 2\}$ und den Erwartungswert der Messung
  $f(x) = x^2$.
\end{exercise}
\begin{solution}
  $\Prob[S] = \pw(1) + \pw(2) = \tfrac23$ und
  $\Ex[f] = \tfrac13\,(1 + 4 + 9) = \tfrac{14}{3}$.
\end{solution}

\begin{tpctable}{c01.tab.sprache}
  \caption{Dieselben Dinge in zwei Sprachen. Links die Wörter der
    Wahrscheinlichkeitstheorie, rechts die der Physik, die dieser Kurs benutzt.}
  \begin{tabular}{lll}
    \toprule
    Symbol & Wahrscheinlichkeitstheorie & Physik \\
    \midrule
    $\Outcomes$ & Stichprobenraum & Ergebnismenge \\
    $\pw$, $\Prob$ & Wahrscheinlichkeitsverteilung, -maß & Präparation \\
    $f$ & Zufallsvariable & Messung \\
    $\ind{S}$ & Indikatorfunktion des Ereignisses $S$ & Ja-Nein-Messung \\
    $\Ex[f]$ & Erwartungswert & Erwartungswert \\
    \bottomrule
  \end{tabular}
\end{tpctable}

\section{Varianz und Determinismus}\label{c01.s3}

Der Erwartungswert sagt, wo die Messwerte im Mittel liegen. Die Varianz sagt, wie weit sie
um diesen Wert streuen. In der Quantenmechanik spielt sie eine große Rolle; bei der
Unschärferelation wird sie uns wieder begegnen.

\tpcdef{D56}

In der Definition ist $f - \Ex[f]$ die Messung mit den Werten $f(x) - \Ex[f]$, die
Abweichung vom Erwartungswert, und $(f - \Ex[f])^2$ ihr Quadrat.

\begin{proposition}[Eigenschaften der Varianz]\label{c01.prop.varianz}\claim[proof-included]{c01.cl.varianz}{derived}
  Sei $\pw$ ein Gewicht auf einer endlichen Ergebnismenge $\Outcomes$, $f$ eine Messung
  und $\alpha$, $\gamma$ reelle Zahlen. Dann gilt:
  \begin{enumerate}
    \item $\Var[f] = \Ex[f^2] - \Ex[f]^2$.
    \item $\Var[f] \geq 0$.
    \item $\Var[f] = 0$ gilt genau dann, wenn $f(x) = \Ex[f]$ für alle Ergebnisse $x$ mit
      $\pw(x) > 0$ ist.
    \item $\Var[\alpha f + \gamma] = \alpha^2\, \Var[f]$.
  \end{enumerate}
\end{proposition}
\begin{proof}
  Wir schreiben $m = \Ex[f]$; das ist eine Zahl. Zu~1: Es ist
  $(f - m)^2 = f^2 - 2 m f + m^2$. Mit der Linearität und der Normierung aus
  Proposition~\ref{c01.prop.erwartung} folgt
  \begin{equation}\label{c01.eq.varianzformel}
    \Var[f] = \Ex[f^2] - 2m\, \Ex[f] + m^2 = \Ex[f^2] - 2m^2 + m^2 = \Ex[f^2] - m^2 .
  \end{equation}
  Zu~2 und~3: Nach Definition ist $\Var[f] = \sum_x \pw(x)\,(f(x) - m)^2$. Jeder Summand
  ist $\geq 0$, also auch die Summe. Sie ist genau dann $0$, wenn jeder Summand $0$ ist,
  und der Summand zu $x$ ist genau dann $0$, wenn $\pw(x) = 0$ oder $f(x) = m$ ist.
  Zu~4: Nach der Linearität ist $\Ex[\alpha f + \gamma] = \alpha m + \gamma$. Also ist
  $\alpha f + \gamma - \Ex[\alpha f + \gamma] = \alpha (f - m)$, das Quadrat davon ist
  $\alpha^2 (f - m)^2$, und sein Erwartungswert ist $\alpha^2\, \Var[f]$.
\end{proof}

\begin{example}[Varianz der Augenzahl]\label{c01.ex.wuerfel-varianz}\claim{c01.cl.wuerfel-varianz}{derived}
  Für den fairen Würfel und $f(x) = x$ ist
  $\Ex[f^2] = \tfrac16\,(1 + 4 + 9 + 16 + 25 + 36) = \tfrac{91}{6}$ und mit
  Gleichung~\eqref{c01.eq.augenzahl-fair}
  \begin{equation}\label{c01.eq.varianz-fair}
    \Var[f] = \tfrac{91}{6} - \bigl(\tfrac72\bigr)^2 = \tfrac{182 - 147}{12} = \tfrac{35}{12} .
  \end{equation}
  Für den gezinkten Würfel aus Gleichung~\eqref{c01.eq.gezinkt} ist
  \begin{equation*}
    \Ex[f^2] = \tfrac{1}{12}\,(1 + 2\cdot 4 + 2\cdot 9 + 2\cdot 16 + 2\cdot 25 + 3\cdot 36)
    = \tfrac{217}{12} ,
  \end{equation*}
  und mit Gleichung~\eqref{c01.eq.augenzahl-gezinkt}
  \begin{equation}\label{c01.eq.varianz-gezinkt}
    \Var[f] = \tfrac{217}{12} - \bigl(\tfrac{47}{12}\bigr)^2
            = \tfrac{2604 - 2209}{144} = \tfrac{395}{144} .
  \end{equation}
\end{example}

\begin{example}[Zwei Ergebnisse]\label{c01.ex.muenze}\claim{c01.cl.muenze}{derived}
  Sei $\Outcomes = \{0, 1\}$ mit $\pw(1) = q$ und $\pw(0) = 1 - q$ für eine Zahl $q$
  zwischen $0$ und $1$, und sei $f(x) = x$. Dann ist $\Ex[f] = q$, und wegen $f^2 = f$ ist
  auch $\Ex[f^2] = q$. Also ist
  \begin{equation}\label{c01.eq.zwei-ergebnisse}
    \Var[f] = q - q^2 = q\,(1 - q) = \tfrac14 - \bigl(q - \tfrac12\bigr)^2 .
  \end{equation}
  Die Varianz ist höchstens $\tfrac14$, und sie ist genau für $q = 0$ und $q = 1$ gleich
  $0$: dann, wenn das Ergebnis feststeht.
\end{example}

Dieselbe Rechnung gilt für jede Ja-Nein-Messung. Nach
Gleichung~\eqref{c01.eq.indikator-regeln} ist $\ind{S}^2 = \ind{S}$, also
\begin{equation}\label{c01.eq.var-indikator}\claim{c01.cl.var-indikator}{derived}
  \Var[\ind{S}] = \Prob[S] - \Prob[S]^2 = \Prob[S]\,\bigl(1 - \Prob[S]\bigr) .
\end{equation}

\begin{tpcfigure}{c01.fig.varianz}
  \caption{(a), (b) Zwei Gewichte mit demselben Erwartungswert $\tfrac72$ der Augenzahl
    $f(x) = x$ (gestrichelte Linie) und verschiedener Varianz: der faire Würfel aus
    Gleichung~\eqref{c01.eq.varianz-fair} und ein Gewicht, das je zur Hälfte auf $3$ und
    $4$ liegt; für dieses ist $\Ex[f^2] = \tfrac{9 + 16}{2}$ und
    $\Var[f] = \tfrac{25}{2} - \tfrac{49}{4} = \tfrac14$. (c) Die Varianz
    einer Ja-Nein-Messung als Funktion von $q = \Prob[S]$ nach
    Gleichung~\eqref{c01.eq.var-indikator}: Sie verschwindet genau an den beiden Enden, wo
    die Antwort feststeht, und ist bei $q = \tfrac12$ am größten.}
\end{tpcfigure}

\subsection{Punktgewichte}\label{c01.s3.1}

Für ein Ergebnis $a \in \Outcomes$ heißt
\begin{equation}\label{c01.eq.punktgewicht}\claim{c01.cl.punktgewicht-def}{source-backed}
  \pw_a(x) =
  \begin{cases}
    1, & x = a, \\
    0, & \text{\tpcifenglish{otherwise}{sonst}}
  \end{cases}
\end{equation}
das \emph{Punktgewicht} auf $a$. Es ist ein Gewicht, und es ist auf einen Punkt
konzentriert: Präpariert wird mit Sicherheit das Ergebnis $a$. Für jede Messung $f$ ist
$\Ex[f] = f(a)$. Jedes Gewicht ist eine Mischung von Punktgewichten, und zwar auf genau
eine Weise:
\begin{equation}\label{c01.eq.zerlegung}\claim{c01.cl.zerlegung}{derived}
  \pw = \sum_{a \in \Outcomes} \pw(a)\, \pw_a .
\end{equation}
Zum Nachprüfen wertet man beide Seiten an einer Stelle $x$ aus: Rechts bleibt nur der
Summand mit $a = x$, und der ist $\pw(x)$. Eindeutig ist die Zerlegung, weil aus
$\pw = \sum_a c_a\, \pw_a$ durch Auswerten an der Stelle $x$ sofort $c_x = \pw(x)$ folgt.

\begin{proposition}[Determinismus ist ein Sonderfall]\label{c01.prop.punktgewicht}\claim[proof-included]{c01.cl.punktgewicht}{derived}
  Für ein Gewicht $\pw$ auf einer endlichen Ergebnismenge $\Outcomes$ sind gleichwertig:
  \begin{enumerate}
    \item $\pw$ ist ein Punktgewicht.
    \item Jedes Ereignis hat die Wahrscheinlichkeit $0$ oder $1$.
    \item Jede Messung hat die Varianz $0$.
    \item $\pw$ ist keine echte Mischung: Ist $\pw = \lambda q + (1 - \lambda) r$ mit
      Gewichten $q$, $r$ und $0 < \lambda < 1$, so ist $q = r = \pw$.
  \end{enumerate}
\end{proposition}
\begin{proof}
  Aus~1 folgt~2: Für $\pw = \pw_a$ ist $\Prob[S] = \ind{S}(a)$, also $0$ oder $1$.

  Aus~2 folgt~1: Für jedes $x$ ist $\pw(x) = \Prob[\{x\}]$ gleich $0$ oder $1$. Die Summe
  aller $\pw(x)$ ist $1$; also ist genau ein $\pw(a)$ gleich $1$ und alle anderen sind $0$.

  Aus~1 folgt~3: Für $\pw = \pw_a$ ist $\Ex[f] = f(a)$, und $a$ ist das einzige Ergebnis
  mit $\pw(x) > 0$. Nach Proposition~\ref{c01.prop.varianz}, Teil~3, ist $\Var[f] = 0$.

  Aus~3 folgt~2: Nach Gleichung~\eqref{c01.eq.var-indikator} ist
  $\Prob[S]\,(1 - \Prob[S]) = \Var[\ind{S}] = 0$, also $\Prob[S] = 0$ oder $\Prob[S] = 1$.

  Aus~1 folgt~4: Sei $\pw_a = \lambda q + (1 - \lambda) r$. Für $x \neq a$ ist
  $0 = \lambda q(x) + (1 - \lambda) r(x)$. Beide Summanden sind $\geq 0$ und
  $\lambda$, $1 - \lambda$ sind positiv, also ist $q(x) = r(x) = 0$. Weil $q$ und $r$
  Gewichte sind, bleibt $q(a) = r(a) = 1$.

  Aus~4 folgt~1: Wir zeigen, dass ein Gewicht, das kein Punktgewicht ist, eine echte
  Mischung ist. Ist $\pw$ kein Punktgewicht, so gibt es zwei verschiedene Ergebnisse $a$
  und $b$ mit $\pw(a) > 0$ und $\pw(b) > 0$. Sei $\varepsilon$ die kleinere der beiden
  Zahlen. Dann sind
  \begin{equation}\label{c01.eq.echte-mischung}
    q = \pw + \varepsilon\,(\pw_a - \pw_b),
    \qquad
    r = \pw - \varepsilon\,(\pw_a - \pw_b)
  \end{equation}
  Gewichte: Die Summe der Werte ist jeweils $1 + \varepsilon\,(1 - 1) = 1$, und negativ
  könnten nur $q(b) = \pw(b) - \varepsilon$ und $r(a) = \pw(a) - \varepsilon$ werden, die
  nach Wahl von $\varepsilon$ aber $\geq 0$ sind. Es ist $q \neq r$, weil
  $q(a) - r(a) = 2\varepsilon > 0$ ist, und $\pw = \tfrac12 q + \tfrac12 r$.
\end{proof}

\begin{tpcfigure}{c01.fig.mischung}
  \caption{Die Konstruktion aus Gleichung~\eqref{c01.eq.echte-mischung} für den gezinkten
    Würfel $\pw$ aus Gleichung~\eqref{c01.eq.gezinkt}, mit $a = 6$, $b = 1$ und
    $\varepsilon = \tfrac{1}{12}$. Von $\pw$ aus (links) wird das Gewicht $\tfrac{1}{12}$
    einmal von $1$ nach $6$ verschoben (Mitte, $q$) und einmal von $6$ nach $1$ (rechts,
    $r$); die veränderten Balken sind schraffiert. Hier ist $r$ der faire Würfel. Der gezinkte Würfel
    ist die Mischung $\pw = \tfrac12 q + \tfrac12 r$ und deshalb kein Punktgewicht.}
\end{tpcfigure}

Im endlichen Modell nennen wir eine Präparation deterministisch, wenn ihr Gewicht ein
Punktgewicht ist. Dann steht das Ergebnis jeder Messung fest: $f$ ergibt mit
Wahrscheinlichkeit $1$ den Wert $f(a)$. So sehen die definitiven Vorhersagen aus, von denen
in Kapitel~0 die Rede war, und so passen sie in eine probabilistische Theorie: als
Präparationen, die keine echten Mischungen sind.
% GAP-TJO: Dass Punktgewichte das Beispiel zu „deterministic as a subset of probabilistic“
% (BRIEF.md:45) sind, ist die Lesart des Kapitelautors; die Proposition sagt nichts über die
% Dynamik deterministischer Theorien (G-0023 d).

\begin{merksatz}[Merksatz]
  Ein Gewicht ist genau dann ein Punktgewicht, wenn jede Messung die Varianz $0$ hat, und
  genau dann, wenn es keine echte Mischung anderer Gewichte ist. Jedes andere Gewicht ist
  auf genau eine Weise eine Mischung von Punktgewichten.
\end{merksatz}

Kapitel~2 nennt Zustände, die keine echten Mischungen sind, \emph{rein}. Die reinen
Zustände dieses endlichen klassischen Modells sind also die Punktgewichte.
Gleichung~\eqref{c01.eq.zerlegung} zerlegt jedes Gewicht eindeutig in Punktgewichte.
Dasselbe Gewicht lässt sich auch herstellen, indem man mit Wahrscheinlichkeit $\pw(a)$ das
Ergebnis $a$ sicher präpariert.
Geometrisch bilden die Gewichte auf $n$ Ergebnissen einen Körper mit $n$ Ecken, den
Punktgewichten: eine Strecke, ein Dreieck, ein Tetraeder\zoomref{c01.v1}.
% GAP-TJO: Die Lesart der Zerlegung als fehlende Information über ein sicher präpariertes
% Ergebnis (nach REVIEW-lec1 aus dem Text genommen) verbindet Gleichung c01.eq.zerlegung mit vl01.tex:83 („Wahrscheinlichkeiten treten
% dort allein wegen fehlender Information auf“); TJO bestätigt oder verwirft sie (G-0023 d).

\begin{vertiefung}[id=c01.v1,
    titel={Die Menge aller Präparationen ist ein Simplex},
    wozu={Der Kern zeigt, dass jedes Gewicht eindeutig aus Punktgewichten gemischt ist.
      Hier wird daraus ein Bild: Für $n$ Ergebnisse bilden die Gewichte eine Strecke, ein
      Dreieck, ein Tetraeder und so weiter, und die Punktgewichte sind die Ecken.},
    voraussetzung={Abschnitt~\ref{c01.s3.1}; Vektoren im $\mathbb{R}^n$. Der Rest dieses
      Kapitels braucht den Block nicht; Kapitel~2 greift das Bild auf.},
    zurueck={Die Gewichte auf $n$ Ergebnissen bilden einen Körper mit $n$ Ecken, den
      Punktgewichten; jeder Punkt darin ist auf genau eine Weise aus den Ecken gemischt.}]
  Sei $\Outcomes = \{1, \dots, n\}$. Ein Gewicht $\pw$ ist dann die Liste
  $(\pw(1), \dots, \pw(n))$ seiner Werte, ein Vektor im $\mathbb{R}^n$. Das Punktgewicht
  $\pw_a$ ist der Vektor mit einer $1$ an der Stelle $a$ und Nullen sonst.

  \begin{proposition}[Die Gewichte bilden ein Simplex]\label{c01.prop.simplex}\claim[proof-included]{c01.cl.simplex}{derived}
    Die Menge aller Gewichte auf $\{1, \dots, n\}$ ist die Menge aller Mischungen
    $\lambda_1 \pw_1 + \dots + \lambda_n \pw_n$ der $n$ Punktgewichte mit
    $\lambda_a \geq 0$ und $\lambda_1 + \dots + \lambda_n = 1$. Die Punktgewichte sind
    affin unabhängig: Aus
    \begin{equation}\label{c01.eq.affin-unabhaengig}
      \alpha_1 \pw_1 + \dots + \alpha_n \pw_n = 0
      \qquad \text{\tpcifenglish{and}{und}} \qquad
      \alpha_1 + \dots + \alpha_n = 0
    \end{equation}
    folgt $\alpha_1 = \dots = \alpha_n = 0$.
  \end{proposition}
  \begin{proof}
    Jede solche Mischung ist ein Gewicht, mit denselben Gründen wie bei
    Gleichung~\eqref{c01.eq.mischung}. Umgekehrt ist jedes Gewicht nach
    Gleichung~\eqref{c01.eq.zerlegung} eine solche Mischung, mit $\lambda_a = \pw(a)$. Zur
    affinen Unabhängigkeit: Die linke Seite der ersten Gleichung in
    \eqref{c01.eq.affin-unabhaengig} ist der Vektor $(\alpha_1, \dots, \alpha_n)$; ist er
    $0$, so sind alle $\alpha_a = 0$. Die zweite Gleichung wird dafür nicht gebraucht.
  \end{proof}

  Die konvexe Hülle von $n$ affin unabhängigen Punkten heißt in Kapitel~2 ein Simplex, und
  eine Theorie, deren Zustandsraum ein Simplex ist, heißt dort klassisch. Für $n = 2$ ist
  das Simplex eine Strecke, für $n = 3$ ein Dreieck (Abbildung~\ref{c01.fig.simplex}), für
  $n = 4$ ein Tetraeder. Die Zahlen $\pw(1), \dots, \pw(n)$ sind zugleich die
  Wahrscheinlichkeiten der Ergebnisse und die Mischungsgewichte der Ecken.

  \begin{tpcfigure}{c01.fig.simplex}
    \caption{Die Menge aller Gewichte. (a) Zwei Ergebnisse: die Strecke zwischen den
      Punktgewichten $\pw_1 = (1,0)$ und $\pw_2 = (0,1)$ in der Ebene der Koordinaten
      $\pw(1)$, $\pw(2)$; eingetragen ist $\pw = (\tfrac14, \tfrac34)$. (b) Drei
      Ergebnisse: ein Dreieck mit den Punktgewichten als Ecken. Eingetragen sind das
      gleichmäßige Gewicht $(\tfrac13, \tfrac13, \tfrac13)$ im Schwerpunkt und das Gewicht
      $(\tfrac12, \tfrac13, \tfrac16)$, das näher an $\pw_1$ liegt. Gefüllte Punkte sind
      Punktgewichte, offene Kreise echte Mischungen.}
  \end{tpcfigure}
\end{vertiefung}

\begin{vertiefung}[id=c01.v2,
    titel={Wie nah liegen relative Häufigkeiten an den Wahrscheinlichkeiten?},
    wozu={Der Kern beschränkt die statistische Lesart auf unabhängige Wiederholungen
      derselben Präparation. Hier wird eine obere Schranke für die Wahrscheinlichkeit
      einer Abweichung der relativen Häufigkeit um mindestens $\varepsilon > 0$ bewiesen;
      diese Schranke nimmt mit der positiven ganzen Anzahl der Wiederholungen ab.},
    voraussetzung={Abschnitte~\ref{c01.s2} und~\ref{c01.s3} bis
      Gleichung~\eqref{c01.eq.var-indikator}. Der Rest des Kurses benutzt nur das Ergebnis
      in Worten.},
    zurueck={Die Wahrscheinlichkeit, dass die relative Häufigkeit eines Ereignisses nach
      $N$ unabhängigen Wiederholungen um mindestens $\varepsilon$ von seiner
      Wahrscheinlichkeit abweicht, ist höchstens $1/(4N\varepsilon^2)$.}]
  Der erste Schritt gilt für jedes Gewicht und jede Messung.

  \begin{lemma}[Ungleichung von Tschebyscheff]\label{c01.lem.tschebyscheff}\claim[proof-included]{c01.cl.tschebyscheff}{derived}
    Sei $\pw$ ein Gewicht auf einer endlichen Ergebnismenge $\Outcomes$, $f$ eine Messung
    und $\varepsilon > 0$. Sei $S$ das Ereignis, dass $f$ um mindestens $\varepsilon$ von
    ihrem Erwartungswert abweicht,
    $S = \{x \in \Outcomes \colon \abs{f(x) - \Ex[f]} \geq \varepsilon\}$. Dann ist
    \begin{equation}\label{c01.eq.tschebyscheff}
      \Prob[S] \leq \frac{\Var[f]}{\varepsilon^2} .
    \end{equation}
  \end{lemma}
  \begin{proof}
    Für $x \in S$ ist $(f(x) - \Ex[f])^2 \geq \varepsilon^2$, und für alle $x$ ist
    $(f(x) - \Ex[f])^2 \geq 0$. Also gilt $(f - \Ex[f])^2 \geq \varepsilon^2\, \ind{S}$ an
    jeder Stelle. Mit Monotonie und Linearität (Proposition~\ref{c01.prop.erwartung}) und
    Gleichung~\eqref{c01.eq.indikator} folgt
    $\Var[f] \geq \varepsilon^2\, \Ex[\ind{S}] = \varepsilon^2\, \Prob[S]$.
  \end{proof}

  Nun zum Modell. Sei $N$ eine positive ganze Zahl und $\Outcomes$ endlich. Ein Ergebnis
  von $N$ Wiederholungen ist eine Liste $(x_1, \dots, x_N)$
  von Ergebnissen aus $\Outcomes$. Wir \emph{nehmen an}, dass die Wiederholungen mit
  derselben Präparation $\pw$ und unabhängig voneinander ausgeführt werden, und
  beschreiben das durch das \emph{Produktgewicht}
  $\pw^{(N)}(x_1, \dots, x_N) = \pw(x_1)\cdots\pw(x_N)$ auf der Menge aller solchen Listen.
  Erwartungswert, Wahrscheinlichkeit und Varianz in diesem Gewicht schreiben wir $\Ex_N$,
  $\Prob_N$ und $\Var_N$.

  \begin{lemma}[Produktgewicht]\label{c01.lem.produkt}\claim[proof-included]{c01.cl.produkt}{derived}
    Das Produktgewicht $\pw^{(N)}$ ist ein Gewicht. Sind $g$ und $w$ Messungen auf
    $\Outcomes$ und $i \neq j$ zwei der Nummern $1, \dots, N$, so gilt für die Messungen
    $(x_1, \dots, x_N) \mapsto g(x_i)$ und $(x_1, \dots, x_N) \mapsto g(x_i)\,w(x_j)$ auf
    den Listen
    \begin{equation}\label{c01.eq.produkt}
      \Ex_N[g(x_i)] = \Ex[g],
      \qquad
      \Ex_N[g(x_i)\,w(x_j)] = \Ex[g]\;\Ex[w] .
    \end{equation}
  \end{lemma}
  \begin{proof}
    Eine Summe über alle Listen von Produkten $u_1(x_1)\cdots u_N(x_N)$ ist das Produkt der
    Einzelsummen,
    \begin{equation*}
      \sum_{x_1, \dots, x_N} u_1(x_1)\cdots u_N(x_N)
      = \Bigl(\sum_{x_1} u_1(x_1)\Bigr)\cdots\Bigl(\sum_{x_N} u_N(x_N)\Bigr) ;
    \end{equation*}
    das ist das Ausmultiplizieren der rechten Seite. Mit $u_l = \pw$ für alle $l$ ergibt das die Summe
    $1$, und alle Werte von $\pw^{(N)}$ sind $\geq 0$. Mit $u_i = \pw\, g$ und $u_l = \pw$
    sonst ist die Einzelsumme zu $i$ gleich $\Ex[g]$, alle anderen sind $1$. Mit
    $u_i = \pw\, g$, $u_j = \pw\, w$ und $u_l = \pw$ sonst sind die Einzelsummen zu $i$ und
    $j$ gleich $\Ex[g]$ und $\Ex[w]$, alle anderen $1$.
  \end{proof}

  \begin{theorem}[Relative Häufigkeiten bei unabhängigen Wiederholungen]\label{c01.thm.haeufigkeit}\claim[proof-included]{c01.cl.gesetz}{derived}
    Sei $\pw$ ein Gewicht auf einer endlichen Ergebnismenge $\Outcomes$, $S$ ein Ereignis
    und $N$ eine positive ganze Zahl. Sei $H_N$ die relative Häufigkeit von $S$ in einer Liste,
    \begin{equation*}
      H_N(x_1, \dots, x_N) = \tfrac{1}{N}\bigl(\ind{S}(x_1) + \dots + \ind{S}(x_N)\bigr) .
    \end{equation*}
    Dann gilt im Produktgewicht $\pw^{(N)}$
    \begin{equation}\label{c01.eq.gesetz-momente}
      \Ex_N[H_N] = \Prob[S],
      \qquad
      \Var_N[H_N] = \frac{\Prob[S]\,(1 - \Prob[S])}{N} ,
    \end{equation}
    und für jedes $\varepsilon > 0$
    \begin{equation}\label{c01.eq.gesetz}
      \Prob_N\bigl[\,\abs{H_N - \Prob[S]} \geq \varepsilon\,\bigr]
      \;\leq\; \frac{\Prob[S]\,(1 - \Prob[S])}{N \varepsilon^2}
      \;\leq\; \frac{1}{4 N \varepsilon^2} .
    \end{equation}
  \end{theorem}
  \begin{proof}
    Nach Gleichung~\eqref{c01.eq.produkt} mit $g = \ind{S}$ ist
    $\Ex_N[\ind{S}(x_i)] = \Prob[S]$ für jedes $i$, und mit der Linearität folgt
    $\Ex_N[H_N] = \tfrac{1}{N}\cdot N\, \Prob[S] = \Prob[S]$. Für die Varianz setzen wir
    $g = \ind{S} - \Prob[S]$; dann ist $\Ex[g] = 0$ und $\Ex[g^2] = \Var[\ind{S}]$. Es ist
    $H_N - \Prob[S] = \tfrac{1}{N}\sum_i g(x_i)$ und
    \begin{equation}\label{c01.eq.gesetz-rechnung}
      \begin{aligned}
        \Var_N[H_N]
        &= \frac{1}{N^2} \sum_{i} \Ex_N[g(x_i)^2]
          + \frac{1}{N^2} \sum_{i \neq j} \Ex_N[g(x_i)\,g(x_j)] \\
        &= \frac{1}{N^2}\cdot N\, \Var[\ind{S}] + 0 .
      \end{aligned}
    \end{equation}
    Nach Gleichung~\eqref{c01.eq.produkt} ist jedes gemischte Glied gleich
    $\Ex[g]\,\Ex[g]$, also $0$. Mit Gleichung~\eqref{c01.eq.var-indikator} ist das die
    zweite Aussage in \eqref{c01.eq.gesetz-momente}. Die erste Ungleichung in
    \eqref{c01.eq.gesetz} ist Lemma~\ref{c01.lem.tschebyscheff} für die Messung $H_N$ im
    Gewicht $\pw^{(N)}$, die zweite folgt aus $q\,(1-q) \leq \tfrac14$
    (Gleichung~\eqref{c01.eq.zwei-ergebnisse}).
  \end{proof}

  Die Schranke fällt wie $1/N$. Sie sagt nicht, dass $h_N$ in jeder Messreihe gegen
  $\pw$ strebt, und sie sagt nichts über eine einzelne Messreihe. Sie sagt, dass
  Messreihen mit großer Abweichung im Modell wenig Gewicht haben. Sie ist außerdem grob:
  Die tatsächliche Wahrscheinlichkeit der Abweichung kann viel kleiner sein als die
  Schranke.

  Für eine Münze lässt sich das Gewicht von $H_N$ ganz angeben. Sei
  $\Outcomes = \{0, 1\}$, $\pw(1) = q$ mit $0 \leq q \leq 1$ und $S = \{1\}$; $N$ sei eine
  positive ganze Zahl. Wir benutzen die Konvention $t^0 = 1$ auch für $t = 0$: Ein
  Produkt ohne Faktoren ist $1$. Eine Liste mit genau $k$ Einsen hat
  das Gewicht $q^k (1-q)^{N-k}$, und es gibt $\binom{N}{k}$ solche Listen, so viele, wie
  es Möglichkeiten gibt, die $k$ Plätze der Einsen unter den $N$ Plätzen auszuwählen. Also
  ist
  \begin{equation}\label{c01.eq.binomial}\claim{c01.cl.binomial}{derived}
    \Prob_N\bigl[H_N = \tfrac{k}{N}\bigr] = \binom{N}{k}\, q^k\, (1-q)^{N-k},
    \qquad k = 0, 1, \dots, N .
  \end{equation}
  Abbildung~\ref{c01.fig.konzentration} zeigt diese Wahrscheinlichkeiten für die faire
  Münze.

  \begin{tpcfigure}{c01.fig.konzentration}
    \caption{Das Gewicht der relativen Häufigkeit $H_N$ von „Zahl“ bei einer fairen Münze,
      $q = \tfrac12$, nach Gleichung~\eqref{c01.eq.binomial} für $N = 4$, $16$ und $64$
      Würfe. Waagerecht steht der Wert $k/N$ von $H_N$; die Höhe eines Strichs ist seine
      Wahrscheinlichkeit, in jedem Bild in Einheiten des größten Werts. Die Klammer
      markiert die Werte mit $\abs{H_N - \tfrac12} < \tfrac14$. Mit wachsendem $N$
      sammelt sich das Gewicht in der Mitte.}
  \end{tpcfigure}
\end{vertiefung}

\section{Kontinuierliche Ergebnisse}\label{c01.s4}

% Stimme: vl01.tex:176-187; v02 [49:00]-[56:35] („wir umgehen das Problem, indem wir die
% wichtigen Eigenschaften nehmen und daraus neue Definitionen machen“; „ich persönlich mag
% diese Notation nicht so gern“).
Bisher war $\Outcomes$ endlich. Für die Quantenmechanik ist aber
$\Outcomes = \mathbb{R}$ entscheidend: Ein Teilchen in einer Dimension hat eine Position,
und die möglichen Positionen sind die reellen Zahlen. Wie man Wahrscheinlichkeiten auf
beliebigen Mengen sauber behandelt, ist Sache der Maßtheorie. Wir umgehen das: Wir nehmen
die wichtigen Eigenschaften des endlichen Falls und machen daraus neue Definitionen, und zwar
nur für den Spezialfall, den wir im Kurs tatsächlich brauchen. In diesem Abschnitt betrachten wir
nur Präparationen auf $\mathbb{R}$, die eine Dichte haben. In diesem Dichtemodell werden
Wahrscheinlichkeiten durch Integration gewonnen, statt Dichtewerte über Punkte zu summieren.
Eine solche Präparation wird durch eine \emph{Wahrscheinlichkeitsdichte} beschrieben:
eine Funktion $\pw$ auf $\mathbb{R}$ mit
\begin{equation}\label{c01.eq.dichte}\claim{c01.cl.dichte}{source-backed}
  \begin{aligned}
    \pw(x) &\geq 0,
    &\int_{-\infty}^{\infty} \pw(x)\,\dd x &= 1, \\
    \Prob[[a,b]] &= \int_a^b \pw(x)\,\dd x,
    &\Ex[f] &= \int_{-\infty}^{\infty} \pw(x)\, f(x)\,\dd x .
  \end{aligned}
\end{equation}
Die beiden Bedingungen der ersten Zeile entsprechen denen an ein Gewicht. Für reelle
Endpunkte $a \leq b$ gibt die zweite Zeile die Intervallwahrscheinlichkeit an. Der
Erwartungswert einer Messung $f$ wird hier verwendet, sofern
$\int_{-\infty}^{\infty} \pw(x)\,\abs{f(x)}\,\dd x < \infty$ ist. Gegenüber dem
endlichen Fall ist also nur die Summe durch ein Integral ersetzt. In der Literatur findet man
für den Erwartungswert auch die Schreibweise $\int \Prob(\dd x)\, f(x)$, mit
$\Prob(\dd x) = \pw(x)\,\dd x$. Ich mag sie nicht besonders und werde sie kaum
verwenden; Sie sollen sie nur erkennen können.

\begin{merksatz}[Eine Dichte ist keine Wahrscheinlichkeit]
  Der Wert $\pw(x)$ einer Dichte ist keine Wahrscheinlichkeit. Wahrscheinlichkeiten sind
  Flächen unter der Kurve, also Integrale von $\pw$ über Mengen. Ein einzelner Punkt hat
  die Wahrscheinlichkeit $0$, und $\pw(x)$ darf größer als $1$ sein.
\end{merksatz}

Beides sieht man an der gleichmäßigen Dichte auf dem Intervall von $0$ bis $\tfrac12$:
\begin{equation}\label{c01.eq.gleich}\claim{c01.cl.gleich}{derived}
  \begin{aligned}
    \pw(x) &=
    \begin{cases}
      2, & 0 \leq x \leq \tfrac12, \\
      0, & \text{\tpcifenglish{otherwise}{sonst}},
    \end{cases}
    &\int_{-\infty}^{\infty} \pw(x)\,\dd x &= 2\cdot\tfrac12 = 1, \\
    \Prob[[a,a]] &= \int_a^a \pw(x)\,\dd x = 0 .
  \end{aligned}
\end{equation}
Die Dichte hat den Wert $2$ und ist trotzdem normiert. Die letzte Gleichung gilt für jede
Dichte und jeden Punkt $a$: Ein Integral über ein Intervall der Länge $0$ ist $0$.

\begin{tpcfigure}{c01.fig.dichte}
  \caption{(a) Die Dichte der Gauß-Verteilung aus Beispiel~\ref{c01.ex.gauss}. Die
    Wahrscheinlichkeit, dass das Ergebnis zwischen reellen Endpunkten $a \leq b$ liegt, ist die schraffierte
    Fläche und nicht die Höhe der Kurve. (b) Die gleichmäßige Dichte aus
    Gleichung~\eqref{c01.eq.gleich}: Sie hat die Höhe $2$, die Fläche darunter ist $1$.}
\end{tpcfigure}

\begin{example}[Gauß-Verteilung]\label{c01.ex.gauss}\claim{c01.cl.gauss}{derived}
  Die Dichte
  \begin{equation}\label{c01.eq.gauss}
    \pw(x) = \frac{1}{\sqrt{2\pi}}\, e^{-\frac12 x^2}
  \end{equation}
  heißt \emph{Normal-} oder \emph{Gauß-Verteilung}, die Glockenkurve. Sie ist eigentlich das
  wichtigste Beispiel der Wahrscheinlichkeitstheorie überhaupt. Sie ist normiert: Die Fläche
  unter der Kurve ist $1$\zoomref{c01.v3}. Wir präparieren ein Teilchen in diesem Zustand
  und schauen am Ende seine Position an. Wie wahrscheinlich ist ein Ergebnis $x \geq 0$?
  Weil $\pw(-x) = \pw(x)$ ist, sind die Flächen links und rechts vom Nullpunkt gleich groß;
  zusammen sind sie $1$. Also ist
  \begin{equation}\label{c01.eq.gauss-halb}
    \int_0^\infty \pw(x)\,\dd x = \int_{-\infty}^0 \pw(x)\,\dd x = \tfrac12 .
  \end{equation}
  Physikalisch gefragt: Ist das Teilchen links oder rechts vom Ursprung? Anschaulich war das
  Ergebnis schon vorher klar. Seien Sie mit dem Argument aber vorsichtig: Es funktioniert
  \emph{nur} für symmetrische Dichten. Für eine unsymmetrische Dichte ist ein solches
  Integral keineswegs offensichtlich.
\end{example}

% Stimme: vl01.tex:204, wörtlich (REVIEW-lec1).
Für Maße auf $\mathbb{R}$ ist es oft besonders einfach, die Erwartungswerte der
Koordinatenfunktion $f(x) = x$ und ihrer Potenzen zu berechnen. Diese Erwartungswerte heißen
\emph{Momente}, sofern die Integrale
$\int_{-\infty}^{\infty} \abs{x}^n\,\pw(x)\,\dd x$ endlich sind:
\begin{equation}\label{c01.eq.momente}\claim{c01.cl.momente}{source-backed}
  M_n = \int_{-\infty}^{\infty} x^n\, \pw(x)\,\dd x ,
  \qquad n = 0, 1, 2, \dots
\end{equation}
Das $n$-te Moment ist also nichts anderes als der Erwartungswert der Messung $x^n$. Es ist
immer $M_0 = 1$: Das ist die Normierung, geometrisch die Fläche unter der Kurve. Die
praktisch wichtigsten Momente sind $M_1$ und $M_2$; $M_1$ heißt \emph{Mittelwert}. Die Varianz der Koordinate ist wie im endlichen Fall
\begin{equation}\label{c01.eq.var-momente}\claim{c01.cl.var-momente}{derived}
  \begin{aligned}
    \Var[\Prob] &= \int_{-\infty}^{\infty} (x - M_1)^2\, \pw(x)\,\dd x \\
    &= M_2 - 2 M_1\cdot M_1 + M_1^2\cdot M_0 = M_2 - M_1^2 ,
  \end{aligned}
\end{equation}
sofern $\int_{-\infty}^{\infty} \abs{x}\,\pw(x)\,\dd x$ und
$\int_{-\infty}^{\infty} x^2\,\pw(x)\,\dd x$ endlich sind. Für die Quantenmechanik ist die Varianz wichtiger als
das zweite Moment; sie kehrt bei der Unschärferelation wieder. Mittelwert und Varianz
tragen dieselbe Information wie $M_1$ und $M_2$.

Die Momente außer $M_0$ müssen nicht existieren.

\begin{example}[Eine Dichte ohne Mittelwert]\label{c01.ex.kein-mittelwert}\claim{c01.cl.kein-mittelwert}{derived}
  Sei $\pw(x) = 1/x^2$ für $x \geq 1$ und $\pw(x) = 0$ sonst. Das ist eine Dichte:
  \begin{equation}\label{c01.eq.kein-mittelwert}
    \begin{aligned}
      \int_1^{b} \frac{\dd x}{x^2} &= 1 - \frac{1}{b} \;\longrightarrow\; 1, \\
      \int_1^{b} x\cdot\frac{\dd x}{x^2} &= \ln b \;\longrightarrow\; \infty
      \qquad (b \to \infty) .
    \end{aligned}
  \end{equation}
  Das Integral für $M_1$ divergiert: Diese Präparation hat keinen Mittelwert, und erst
  recht keine Varianz.
\end{example}

\begin{vertiefung}[id=c01.v3,
    titel={Normierung und Momente der Gauß-Verteilung},
    wozu={Beispiel~\ref{c01.ex.gauss} behauptet, die Gauß-Dichte sei normiert. Hier wird
      das nachgerechnet, zusammen mit Mittelwert und Varianz.},
    voraussetzung={Abschnitt~\ref{c01.s4}; Integrale in zwei Variablen, Polarkoordinaten,
      partielle Integration. Der Kern benutzt nur das Ergebnis.},
    zurueck={Die Gauß-Dichte $e^{-x^2/2}/\sqrt{2\pi}$ ist normiert, ihr Mittelwert ist $0$
      und ihre Varianz $1$.}]
  Sei $I = \int_{-\infty}^{\infty} e^{-x^2/2}\,\dd x$. Das Integral existiert, weil
  $e^{-x^2/2} \leq e^{-\abs{x}/2}$ für $\abs{x} \geq 1$ ist. Man berechnet $I^2$ als
  Integral über die Ebene und wechselt zu Polarkoordinaten $x = r\cos\varphi$,
  $y = r\sin\varphi$ mit $\dd x\,\dd y = r\,\dd r\,\dd\varphi$:
  \begin{equation}\label{c01.eq.gauss-norm}\claim{c01.cl.gauss-norm}{derived}
    \begin{aligned}
      I^2 &= \int_{-\infty}^{\infty}\!\int_{-\infty}^{\infty} e^{-(x^2 + y^2)/2}\,\dd x\,\dd y
           = \int_0^{2\pi}\!\int_0^\infty e^{-r^2/2}\, r\,\dd r\,\dd\varphi \\
          &= 2\pi\,\Bigl[-e^{-r^2/2}\Bigr]_0^\infty = 2\pi .
    \end{aligned}
  \end{equation}
  Weil $I > 0$ ist, folgt $I = \sqrt{2\pi}$, und die Dichte \eqref{c01.eq.gauss} hat das
  Integral $1$.

  % Stimme: v02 [63:27]-[65:07] („intuitiv soll klar sein, dass die erwartete Position null
  % ist; bewiesen wird es trotzdem“).
  Dass die erwartete Position null ist, ist intuitiv klar; bewiesen wird es trotzdem. Für die
  Momente gilt
  \begin{equation}\label{c01.eq.gauss-momente}\claim{c01.cl.gauss-momente}{derived}
    \begin{aligned}
      M_1 &= 0, \\
      M_2 &= \frac{1}{\sqrt{2\pi}}\int_{-\infty}^{\infty} x\cdot x\,e^{-x^2/2}\,\dd x \\
          &= \frac{1}{\sqrt{2\pi}}\Bigl[-x\,e^{-x^2/2}\Bigr]_{-\infty}^{\infty}
            + \frac{1}{\sqrt{2\pi}}\int_{-\infty}^{\infty} e^{-x^2/2}\,\dd x
          = 0 + 1 .
    \end{aligned}
  \end{equation}
  $M_1$ verschwindet, weil $x\,e^{-x^2/2}$ eine ungerade Funktion ist und ihr Integral
  über $x \geq 0$ existiert: $\int_0^\infty x\,e^{-x^2/2}\,\dd x = 1$. Bei $M_2$ wurde
  partiell integriert, mit $x\,e^{-x^2/2}$ als Ableitung von $-e^{-x^2/2}$; der Randterm
  verschwindet, weil $x\,e^{-x^2/2}$ für $x \to \pm\infty$ gegen $0$ geht. Nach
  Gleichung~\eqref{c01.eq.var-momente} ist die Varianz $M_2 - M_1^2 = 1$.
\end{vertiefung}

Damit steht die Sprache für die Modelle dieses Kapitels bereit: Präparationen sind
Gewichte auf endlichen Ergebnismengen oder die hier betrachteten Dichten auf $\mathbb{R}$,
Messungen sind Funktionen, und beide zusammen ergeben bei endlichen Integralen
Erwartungswerte. Kapitel~2 löst sich von der
Annahme, dass Präparationen Gewichte auf einer Ergebnismenge sind, und behält nur, was
man mit ihnen tun kann: mischen und Ja-Nein-Fragen stellen.
